\chapter{Introduction}

After the acknowledgements and abstract of your dissertation you should have a table of contents. If you wish you can also have a list of figures and a list of tables. Be aware though, that if you have long captions for the figures and tables, these lists can be very long and not necessarily very useful (you can also consider using truncated captions for the lists of figures and tables). The table of contents and the lists of figures and tables do not contribute towards the page count so you may wish to number these pages with Roman numbers. 

\section{Structure}

The first chapter of your dissertation should be an introduction. If you used Roman numbers for the page numbering of the table of contents etc, you should start using  standard Arabic numbers now. Ideally the first page of the introduction should be numbered page 1 as it is the first page that counts towards the page limit.

Think about sectioning your work nicely with sections, subsections and subsubsections.

\subsection{Sectioning}

Be consistent with capitalisation of titles of sections etc. Either capitalise all words (apart from of, the, a etc.) or just capitalise the first word and proper nouns. Choose one of these conventions and stick to it!

\subsubsection{Numbering of Sections}

By default sections and subsections will be numbered in the main body of the text and in the table of contents. You may wish to change this so that subsubsections are also numbered in the main text and perhaps in the table of contents. Commands of the form \verb|\setcounter{secnumdepth}{3}| and \verb|\setcounter{tocdepth}{3}| will do this for you.

\section{About the Introduction}

Use the introduction to explain what the problem is you are going to solve and how you will attack this (methods). You may wish to have a literature review in the introduction or in a separate chapter.

\section{Main Contributions}

The examiners have suggested that a ``main contributions'' section in the introduction is really useful, so please ensure you have such a section! You can use this to explain what you believe are the main results of your work. If any of what you have done is novel, it can be helpful to highlight that here. If you have worked on an industrial project you could state how and why your results are useful to your industrial partner.

\section{Code}

Here you may have a link to your GitHub repository (recommended but not mandatory). You can also make clear whether you have written code from scratch, or based it on other code, or had help from ChatGPT.

\section{Structure of the Dissertation}

It is common, although not compulsory, to end the introduction with a brief section explaining the structure of the dissertation. Often this goes along the lines of:

The remainder of the dissertation is structured as follows. In Chapter 2 we discuss blah blah. In Chapter 3 we blah blah. Finally in Chapter 4 we summarise the work we have done and suggest some areas for future research.

